\documentclass[a4paper,12pt,oneside]{report} %syntonly]
%\usepackage{pstricks,listings} %% Package Tres complet! Une vraie Usine a Gas !!
\usepackage[latin1]{inputenc}     % pour Unix Latin1 
% \usepackage[applemac]{inputenc} % pour utilisateurs mac
% \usepackage[ansinew]{inputenc}  % pour Windows ANSI (yen a encore qui utilisent ca? :s)
\usepackage{verbatim} 
%\usepackage{fullpage}
\usepackage{aeguill}
\usepackage{amsfonts}
\usepackage{xspace}
\usepackage{amsmath}
\usepackage{textcomp} 
\usepackage{cancel} 
\usepackage{mathcomp} 
\usepackage{makeidx}
\usepackage{lscape} % pour avoir une page en paysage
%\usepackage{acronym}
%\usepackage{picins}
%latex RapportKR.tex
%makeindex RapportKR.idx
%latex RapportKR.tex
\usepackage{fancybox}
\usepackage{shapepar}	% pour les mises en forme de texte (losange etc)
\usepackage{fancyhdr}
\usepackage{textcomp}  % pour les caract�res sp�ciaux (copyleft etc)
\usepackage{marvosym} % caractere electrostatic, laserbeam...
\usepackage{supertabular}% tableaux sur plusieurs pages
\usepackage{ifsym} % caract�res scp�ciaux
\usepackage{longtable}
\usepackage{listings}
\usepackage{multicol}
\usepackage{upgreek}
%\usepackage{picins}
%\usepackage{frbib}
%\usepackage{lastpage}
\usepackage[english,french]{babel}
%%%%%%%%%%%%% parallel, multicolpar et parcolumns.%%%%%%%%%%%%%%%%%%%%%% pour des documents multicolones bilingue.
\usepackage[pdftex]{graphicx, color}
\usepackage{pdfpages}
\graphicspath{{img/}} % La ou j'ai mes images
\DeclareGraphicsExtensions{.jpg,.png}
\pdfoutput=1
\usepackage[pdftex,
	a4paper,
        bookmarks = true,           % Signets
        bookmarksnumbered = true,   % Signets numerotes
        pdfpagemode = None,         % Signets/vignettes fermes a l'ouverture
        %pdfstartview = FitH,        % La page prend toute la largeur
	pdfstartview = FitV,        % La page prend toute la hauteur
	colorlinks=true,
	citecolor=black,urlcolor=blue,linkcolor=red,
	pdfauthor={K�vin Raymond},
	pdftitle={Stage au CERN},
 	pdfsubject={Rapport de stage},
	pdfkeywords={Xbee, pt100, PIC, GNU, Linux, CAO, gEDA},
	plainpages=false,
	pdfpagelabels,
	bookmarks,
	breaklinks=true,
   	hyperindex,
	linktocpage=true % pour colorier seulement le num�ros dans la TOC	
]{hyperref}
%
\pdfcompresslevel=9
%
%
%%% style sans serif pour les url
\urlstyle{sf}
%
%%%% debut macro %%%%  Pour mettre en forme Chapitre
\makeatletter
\def\@makechapterhead#1{%
  \vspace*{2\p@}%
  {\parindent \z@ \raggedright \normalfont
    \interlinepenalty\@M
    \ifnum \c@secnumdepth >\m@ne
        \Huge\bfseries \thechapter\quad
    \fi
    \Huge \bfseries #1\par\nobreak
    \vskip 40\p@
  }}
%
\def\@makeschapterhead#1{%
  \vspace*{2\p@}%
  {\parindent \z@ \raggedright
    \normalfont
    \interlinepenalty\@M
    \Huge \bfseries  #1\par\nobreak
    \vskip 40\p@
  }}
\makeatother

%%% Reglage des marges
\setlength{\hoffset}{-25pt}  	
\setlength{\voffset}{-18pt}  	
\setlength{\oddsidemargin}{50pt} 	% Marge gauche sur pages impaires
\setlength{\evensidemargin}{9pt} 	% Marge gauche sur pages paires
\setlength{\marginparwidth}{54pt} 	% Largeur de note dans la marge
\setlength{\textwidth}{420pt} 	% Largeur de la zone de texte (17cm)
\setlength{\marginparsep}{7pt} 	% S�paration de la marge
\setlength{\topmargin}{0pt} 	% Pas de marge en haut
\setlength{\headheight}{33pt} 	% Haut de page
\setlength{\headsep}{40pt} 	% Entre le haut de page et le texte
\setlength{\paperwidth}{600pt}
\setlength{\footskip}{27pt} 	% Bas de page + s�paration
\setlength{\textheight}{650pt} 	% Hauteur de la zone de texte 
%

%%%% debut macro %%%%
\newenvironment{changemargin}[2]{\begin{list}{}{%
\setlength{\topsep}{0pt}%
\setlength{\leftmargin}{0pt}%
\setlength{\rightmargin}{0pt}%
\setlength{\listparindent}{\parindent}%
\setlength{\itemindent}{\parindent}%
\setlength{\parsep}{0pt plus 1pt}%
\addtolength{\leftmargin}{#1}%
\addtolength{\rightmargin}{#2}%
}\item }{\end{list}}
%%%% fin macro %%%%


\pagestyle{fancy} %�viter le style par d�faut.
%
%% Def des entetes, ce referer a http://www.trustonme.net/didactels/250.html
\fancyhead[LO]{\textsc{\rightmark} \rule[0.5ex]{\textwidth}{0.4mm}}
%\fancyhead[C]{}
%\fancyhead[R]{\includegraphics[width=1cm]{./img/AB.jpg}}
\fancyfoot[R]{\footnotesize{\textsc{Rapport de Stage}}}	
\fancyfoot[C]{\footnotesize{- \thepage} -}	
\fancyfoot[L]{\footnotesize{\Laserbeam\ \textit{K�vin Raymond}}}	
%
%%% Ajouter une fine ligne en haut et en bas
%\renewcommand{\headrulewidth}{3pt}
\renewcommand{\footrulewidth}{0.4pt}

%\pagestyle{headings}
%
%
%
%
%
%
%
%
%%%%%%%%% Mise en Forme page de garde %%%%%%%%%%%%%%%%%%%%
%######################## D�finition du titre ######################
\title{
\thispagestyle{empty}
\vspace{-40mm}
   \begin{minipage}[c][160pt][c]{.42\linewidth} %[hauteur][position]{largeur}
	{\normalsize 
		\begin{flushleft}
		\vspace{3mm}
		Enseignement des r�seaux\\
		Universit�e de Bretagne-Sud\\
		Formation ing�nieur cyberd�fense\\
		\url{http://www-ensibs.univ-ubs.fr/}\\
		\vspace{2mm}
		\includegraphics[width=7cm]{./img/ENSIBS_Logo}
		\vspace{3mm}
		\end{flushleft}
	}
   \end{minipage}\hfill
\hspace{15mm}
   \begin{minipage}[c][160pt][c]{.42\linewidth} %[hauteur][position]{largeur}
	{\normalsize 
		\begin{flushright}
		\vspace{3mm}
		Liebhardt J�r�mie\\
		Mollard R�mi\\
		Coueron Adrien\\
		Gambart Thomas\\
		Ann�e 2015\\
		\vspace{30mm}
		\end{flushright}
	}
   \end{minipage}
\vfill
\vspace{30mm}
%
%
   {\Huge \textsc{Rapport de projet}}\\
\vspace{6mm}
%\Ovalbox{
\rule{\linewidth}{0.8mm}\vspace{6mm}\\
    \begin{minipage}[c]{0.70\linewidth} %[hauteur][position]{largeur}
	{\Huge
		Installation et configuration 
		d'un rogue access point
	}
	\vspace{6mm}
   \end{minipage}
\rule{\linewidth}{0.8mm}\vspace{3mm}\\
%}
%
\vspace{5cm}
%
   \begin{minipage}[c]{.46\linewidth}
	{\normalsize 
		\begin{flushleft}
		\vspace{2mm}
		\end{flushleft}
	}
   \end{minipage}\hfill
%
   \begin{minipage}[c]{.46\linewidth}
	{\normalsize 
		\begin{flushright}
		Enseignant\\
		\vspace{2mm}
		\textsc{Ma�l Auzias}\\
		\end{flushright}
	}
   \end{minipage}
\date{} % Je ne veux pas afficher la date
}  %fin du titre

\author{}
%\end{flushright}

%
%
%
%
%
%
%
%
\makeatletter
\renewcommand{\@makefntext}[1]%
{\noindent\llap{\@makefnmark\,}#1}
\renewcommand{\footnoterule}{\vspace*{-3pt}%
\noindent
\hspace*{-1em}
\rule{.4\columnwidth}{0.4pt}\vspace*{2.6pt}}
\makeatother 



%
%%%%%%%%%%%%%%%% Fin des includes... %%%%%%%%%%%%%

\selectlanguage{french}
\makeindex % Au travail !
\begin{document}
\setlength{\headheight}{33pt} 	% Haut de page
\setlength{\headsep}{20pt} 	% Entre le haut de page et le texte
\setlength{\textheight}{950pt} 	% Hauteur de la zone de texte 

%%
\maketitle
\pagestyle{plain}
\renewcommand{\headrulewidth}{0pt}

\setcounter{page}{1} %On definit un compteur a 1
\pagenumbering{roman} %permet une num�rotation en chiffre romain
\thispagestyle{empty} \mbox{ }
\cleardoublepage


\setlength{\headheight}{33pt} 	% Haut de page
\setlength{\headsep}{40pt} 	% Entre le haut de page et le texte
\setlength{\textheight}{650pt} 	% Hauteur de la zone de texte 


%\thispagestyle{headings}
%%\thispagestyle{fancy}

\section*{Remerciements}
Tout d'abord, j'exprime ma profonde gratitude envers Javier SERRANO pour m'avoir pris en stage, ainsi qu'� Philippe NOUCHI sans qui il aurait �t� plus dur d'entrer au CERN. Bien s�r, je les remercie �galement pour leurs nombreux conseils. Ils m'ont permis d'organiser mon travail d'un point de vue professionnel. 

Je me dois �galement de remercier Nicolas DE METZ-NOBLAT, un administrateur qui s'est occup� de la r�installation de mon syst�me et Christine GAYRAUD pour le suivit des diff�rentes commandes.

J'adresse aussi mes remerciements aux autres coll�gues pour la bonne ambiance et les quelques conseils toujours bienvenus. Plus particuli�rement � Matthieu CATTIN et Olivier BARRI\` ERE pour la r�alisation du prototype.

C'est un r�el plaisir de travailler dans une telle �quipe.


\section*{Note}
\diamondpar{Ce projet est r�alis� � base de Pic 18F2550 mais peut facilement �tre adapt� pour d'autres microcontr�leurs. De plus, ce rapport a �t� r�dig� en \LaTeX\ en accord avec le monde du logiciel libre. Ce langage �tant compil�, il permet �galement une mise en page de qualit� et normalis�e. Je vous recommande d'ailleurs la classe Beamer de \LaTeX\ permettant de r�aliser rapidement des pr�sentations et des transparents de tr�s grandes qualit�s. \og La connaissance s'acquiert par l'exp�rience, tout le reste n'est que de l'information\fg\ - Albert Einstein}


 %% Remerciement + note
%%
\renewcommand{\contentsname}{\vspace{-30pt} 
Sommaire} %Pour parler en francais et tout entrer sur 1 seule page
\tableofcontents




%%%%%%%%%%%%%%%%%%%%%%%%%%%%%%%%%%%%%%%%%%%%%%%%%%%%%%%%%%%%%%%%%%%%%%%%%%%%%%%
%\listoffigures   
%\listoftables  
%%%%%%%%%%%%%%%%%%%%%%%%%%%%%%%%%%%%%%%%%%%%%%%%%%%%%%%%%%%%%
\cleardoublepage
\setcounter{page}{1} % On r�initialise le coompteur des pages
\pagenumbering{arabic} % On repasse en num�rotation en chiffre Arabe
%%%%% On redefinit le style fancy{plain} pour les pages sp�ciales (de nouveaux chapitres) : 
\fancypagestyle{plain}{
\fancyhead[LO]{}
\fancyhead[R]{}
\fancyhead[C]{}
}
\pagestyle{fancy}
\renewcommand{\headrulewidth}{3pt}
\chapter*{Pr�sentation}
\addcontentsline{toc}{chapter}{Introduction}
\section*{Nom du projet}
Nous avons d�cid� de nommer notre projet "WifiSansPrix". Le nom scientifique du projet est : Rogue Access Point.
\newline
Notre choix pour le nom du projet est enti�rement tir� du but final. 

\section*{Objectifs du projet}
Le principal objectif de ce projet est d'associer plusieurs technologies de l'informatique et des r�seaux pour mettre en oeuvre une attaque de type "man in the middle" sur un r�seau wifi.
\newline

Ce projet consiste � usurper l'identit� d'une borne wifi pour que le traffic r�seau des utilisateurs connect�s soit automatiquement redirig� vers notre borne wifi corrompue.
\newline
\newline
Afin d'atteindre cet objectif nous allons proc�der ainsi :
\begin{itemize}\renewcommand{\labelitemi}{$\bullet$}
\vspace{3mm}
\item
Cr�ation d'un point d'acc�s � l'aide de l'outil Airbase-ng
\item
Mise en place d'un serveur DHCP pour les clients
\item
Mise en place d'un serveur DNS pour rediriger certaines demandes (Facebook, maBanque, Gmail...) vers un site pi�g�.
\item
Mise en place d'une proc�dure ("Hame�onnage") pour r�cup�rer la cl� wifi de la vraie borne wifi
\end{itemize}
\section*{Nouvelles comp�tences vis�es}
Ce projet a pour but de nous faire comprendre l'installation, la configuration et l'analyse des outils cit�s pr�c�demment.
\newline
\newline
Voici les notions qui serons abord�es pour ce projet :
\begin{itemize}\renewcommand{\labelitemi}{$\bullet$}
\vspace{3mm}
\item
Fonctionnement d'une borne wifi ( Analyse des trames avec wireshark )
\item
Utilisation des outils de la suite aircrack-ng
\item
Analyse du fonctionnement d'un serveur DHCP
\item
Analyse du fonctionnement d'un serveur DNS
\item
D�veloppement de plusieurs sites web (Apache / SQL)
\item
Analyse du SSL Stripping
\end{itemize}
\newpage
\section*{Organisation}
Pour l'organisation de notre travail nous avons pris la d�cision de travailler de la mani�re suivante :
\begin{itemize}\renewcommand{\labelitemi}{$\bullet$}
\vspace{3mm}
\item
Partage des documents, fichiers de configuration, ... via github
\item
Communication par t�l�phone et email
\item
Deux cr�neaux de 2h par semaine pour travailler tous ensembles
\end{itemize}

 %% INTRODUCTION
%\input{./source/objectifs} %% Outils de developpement
%\input{./source/objectifs_attein_comment} %%  ----- TRAVAUX R�ALIS�S
%\input{./source/demo}
%\input{./source/objectifs_non_attein_pk} %% -----
%\input{./source/quid_appri} %% -----
%\input{./source/cmt_tester_legalement} %% IMPRESSIONS PERSONNELLES 
%\input{./source/Conclusion} %% CONCLUSION

\pagestyle{plain}
%\bibliographystyle{unsrt}
%\bibliography{rapport}
%\newpage
%\chapter*{Glossaire}
\addcontentsline{toc}{chapter}{Glossaire}
\begin{changemargin}{-13mm}{-13mm}
 

\begin{multicols}{2}
\setlength{\columnsep}{36pt}
\setlength{\columnseprule}{0.3pt}

\noindent\textbf{Ductile}~Se dit d'un mat�riau capable de se d�former plastiquement sans se rompre.\newline
\textbf{Justesse}~Qualit� du capteur pour donner des indications �gales au mesurande.\newline
\textbf{Kynaptic}~Gestionnaire de paquet pour KDE.\newline
\textbf{Lin�arit�}~Sensibilit� du capteur, ind�pendante du mesurande.\newline
\textbf{Mesurande}~Grandeur physique � mesurer (temp�rature, pressions,~\dots).\newline
\textbf{Fid�lit�}~L'aptitude d'un capteur � donner des indications identiques pour une m�me valeur de mesurande.\newline
\textbf{Finesse}~Aptitude du capteur � fournir une grandeur mesur�e sans la modifier par sa pr�sence.\newline
\textbf{Pr�cision}~Pour un capteur, aptitude � donner une indication proche de la vraie valeur.\newline
\textbf{Rapidit�}~Temps de r�ponse, li�e � la fr�quence propre du capteur.\newline
\textbf{R�solution}~La plus petite valeure de mesurande que le capteur peut d�tecter.\newline
\textbf{Synaptic}~Gestionnaire de paquet sous Linux (Debian), voir aussi Kynaptic.\newline

\section*{Les acronymes}
\noindent\textbf{CAN}~Convertisseur Analogique Num�rique\newline
\textbf{CAO}~Conception Assist� par Ordinateur\newline
\textbf{GPS}~Global Positioning System\newline
\textbf{HID}~Human Interface Device\newline
\textbf{ICD(2)}~In Circuit Debugger (2)\newline
\textbf{JFET}~ Junction Field Effect Transistor\newline
\textbf{PIC}~Nom des Microcontr�leurs con�us par Microchip\newline
\textbf{PCB}~Printed Circuit Board\newline
\textbf{SDCC}~Small Device C Compileur\newline
\textbf{USB}~Universal Serial Bus\newline

\end{multicols}

\end{changemargin}


\chapter*{Bibliographie}
\addcontentsline{toc}{chapter}{Bibliographie}

\noindent\url{http://press.web.cern.ch/}, toutes les nouvelles du CERN, toutes les donn�es et dates importantes.\newline
\url{http://fr.wikipedia.org/}, diverses d�finitions.\newline
\url{http://www.iut-lannion.fr/LEMEN/MPDOC/CMMEST/mtherres.htm}, mesure de temp�rature.\newline
\url{http://aviatechno.free.fr/thermo/rtd03.php},mesure de temp�rature.\newline
\url{http://perso.orange.fr/michel.hubin/physique/elec/chap_aop2.htm}, divers sch�mas �lectroniques. \newline
\url{http://forums.futura-sciences.com/forum40.html}, un forum avec des sujets interr�ssants sur la Pt100.\newline
\url{http://www.planete-sciences.org/forums/}, forum tr�s int�ressant sur l'utilisation du module Xbee. Voir section Robotique.\newline
\url{http://www.jeffreyleblanc.org/energy/xbee.htm}, sources pour les commandes AT du module Xbee.\newline
\url{http://www.todopic.com.ar/foros/index.php?topic=16393.0;topicseen}, un bon exemple sur l'utilisation du module Xbee.\newline
\url{http://www.totofweb.net/}, un beau projet avec le module Xbee.\newline
\url{http://itp.nyu.edu/~raf275/meshnetworking/XBee/XBee_firmware_upgrade.html}, des explications sur la mise � jour du module Xbee ("~upgrading firmware~").\newline
\url{http://pic18fusb.online.fr/}, les PIC pour l'USB en HID.\newline
\url{http://www.analog.com/UploadedFiles/Associated_Docs/49470200sscsect2.PDF}, un dossier tr�s int�ressant en rapport avec le pont de Wheatstone et la Pt100. Dommage que je ne l'ai trouv� qu'� la fin, il m'a permis de v�rifi� mon sch�ma, puisqu'il est reproduit � l'identique (page 11 figure 2.12).\newline
\url{http://lpm.epfl.ch/webdav/site/lpm/shared/enseignement/divers/060914\%20Mesure\%20de\%20temp�rature\%20Pt100-Pt1000.pdf}, un dossier permettant d'ajuster l'erreur de l'ajustement quadratique standard pour les temp�ratures n�gatives de la Pt100 (et Pt1000).\newline
\url{http://fr.wikipedia.org/wiki/Pt100} pour une sonde encore plus pr�cise.\newline
\url{http://aviatechno.free.fr/thermo/rtd03.ph} des explications sur la lin�arit� de la Pt100.\newline



\renewcommand{\headrulewidth}{0pt}
\cleardoublepage
\thispagestyle{empty}

%% Annexes

%\appendix

%\setlength{\oddsidemargin}{10pt} 	% Marge gauche sur pages impaires
%\setlength{\marginparwidth}{10pt} 	% Largeur de note dans la marge

%\addtocontents{toc}{\newpage \vskip 10 mm \Large \centerline {\textsc{Annexes}} \normalsize \vskip 5mm}
%\part*{\Huge{\textsc{Annexes}}}
%\setcounter{page}{1} % R�initialisation du compteur
%\pagenumbering{arabic}
%
%
%\pagestyle{headings}
%\thispagestyle{empty}


\chapter{Tables Pt100} \label{tablePt100}

\newcommand\espace{\vrule height 14pt width 0pt}
%%%

\subsubsection{Table de r�sistivit� de la Pt100 (valeur � multiplier par 10 pour avoir la table de la Pt1000)}


\large{Selon ITS-90/DIN EN 60751}
Unit�s : T en \degres C et R en $\Omega$

\begin{changemargin}{-20mm}{-15mm}
\vspace{4mm}
%\begin{table}
\footnotesize{
\begin{center}
\begin{supertabular}{|c||c|c|c|c|c|c|c|c|c|c||c|}

   	\hline
	\textbf{\degres C} & \textbf{0} & \textbf{1} & \textbf{2} & \textbf{3} & \textbf{4} & \textbf{5} & \textbf{6} & \textbf{7} & \textbf{8} & \textbf{9} & \textbf{\degres C}\\\hline\hline
	-200.00&  18.52& & & & & & & & & &                                              -200.00 \\\hline
	-190.00&  22.83&  22.40&  21.97&  21.54&  21.11&  20.68&  20.25&  19.82&  19.38&  18.95& -190.00 \\\hline
	-180.00&  27.10&  26.67&  26.24&  25.82&  25.39&  24.97&  24.54&  24.11&  23.68&  23.25& -180.00 \\\hline
	-170.00&  31.34&  30.91&  30.49&  30.07&  29.64&  29.22&  28.80&  28.37&  27.95&  27.52& -170.00 \\\hline
	-160.00&  35.54&  35.12&  34.70&  34.28&  33.86&  33.44&  33.02&  32.60&  32.18&  31.76& -160.00 \\\hline
	-150.00&  39.72&  39.31&  38.89&  38.47&  38.05&  37.64&  37.22&  36.80&  36.38&  35.96& -150.00 \\\hline
	-140.00&  43.88&  43.46&  43.05&  42.63&  42.22&  41.80&  41.39&  40.97&  40.56&  40.14& -140.00 \\\hline
	-130.00&  48.00&  47.59&  47.18&  46.77&  46.36&  45.94&  45.53&  45.12&  44.70&  44.29& -130.00 \\\hline
	-120.00&  52.11&  51.70&  51.29&  50.88&  50.47&  50.06&  49.65&  49.24&  48.83&  48.42& -120.00 \\\hline
	-110.00&  56.19&  55.79&  55.38&  54.97&  54.56&  54.15&  53.75&  53.34&  52.93&  52.52& -110.00 \\\hline
	-100.00&  60.26&  59.85&  59.44&  59.04&  58.63&  58.23&  57.82&  57.41&  57.01&  56.60& -100.00 \\\hline
	 -90.00&  64.30&  63.90&  63.49&  63.09&  62.68&  62.28&  61.88&  61.47&  61.07&  60.66&  -90.00 \\\hline
	 -80.00&  68.33&  67.92&  67.52&  67.12&  66.72&  66.31&  65.91&  65.51&  65.11&  64.70&  -80.00 \\\hline
	 -70.00&  72.33&  71.93&  71.53&  71.13&  70.73&  70.33&  69.93&  69.53&  69.13&  68.73&  -70.00 \\\hline
	 -60.00&  76.33&  75.93&  75.53&  75.13&  74.73&  74.33&  73.93&  73.53&  73.13&  72.73&  -60.00 \\\hline
	 -50.00&  80.31&  79.91&  79.51&  79.11&  78.72&  78.32&  77.92&  77.52&  77.12&  76.73&  -50.00 \\\hline
	 -40.00&  84.27&  83.87&  83.48&  83.08&  82.69&  82.29&  81.89&  81.50&  81.10&  80.70&  -40.00 \\\hline
	 -30.00&  88.22&  87.83&  87.43&  87.04&  86.64&  86.25&  85.85&  85.46&  85.06&  84.67&  -30.00 \\\hline
	 -20.00&  92.16&  91.77&  91.37&  90.98&  90.59&  90.19&  89.80&  89.40&  89.01&  88.62&  -20.00 \\\hline
	 -10.00&  96.09&  95.69&  95.30&  94.91&  94.52&  94.12&  93.73&  93.34&  92.95&  92.55&  -10.00 \\\hline
	   0.00& 100.00&  99.61&  99.22&  98.83&  98.44&  98.04&  97.65&  97.26&  96.87&  96.48&    0.00 \\\hline
\hline
	\textbf{\degres C} & \textbf{0} & \textbf{1} & \textbf{2} & \textbf{3} & \textbf{4} & \textbf{5} & \textbf{6} & \textbf{7} & \textbf{8} & \textbf{9} & \textbf{\degres C}\\\hline\hline
	   0.00& 100.00& 100.39& 100.78& 101.17& 101.56& 101.95& 102.34& 102.73& 103.12& 103.51&    0.00 \\\hline
	  10.00& 103.90& 104.29& 104.68& 105.07& 105.46& 105.85& 106.24& 106.63& 107.02& 107.40&   10.00 \\\hline
	  20.00& 107.79& 108.18& 108.57& 108.96& 109.35& 109.73& 110.12& 110.51& 110.90& 111.29&   20.00 \\\hline
	  30.00& 111.67& 112.06& 112.45& 112.83& 113.22& 113.61& 114.00& 114.38& 114.77& 115.15&   30.00 \\\hline
	  40.00& 115.54& 115.93& 116.31& 116.70& 117.08& 117.47& 117.86& 118.24& 118.63& 119.01&   40.00 \\\hline
	  50.00& 119.40& 119.78& 120.17& 120.55& 120.94& 121.32& 121.71& 122.09& 122.47& 122.86&   50.00 \\\hline
	  60.00& 123.24& 123.63& 124.01& 124.39& 124.78& 125.16& 125.54& 125.93& 126.31& 126.69&   60.00 \\\hline
	  70.00& 127.08& 127.46& 127.84& 128.22& 128.61& 128.99& 129.37& 129.75& 130.13& 130.52&   70.00 \\\hline
	  80.00& 130.90& 131.28& 131.66& 132.04& 132.42& 132.80& 133.18& 133.57& 133.95& 134.33&   80.00 \\\hline
	  90.00& 134.71& 135.09& 135.47& 135.85& 136.23& 136.61& 136.99& 137.37& 137.75& 138.13&   90.00 \\\hline
	 100.00& 138.51& 138.88& 139.26& 139.64& 140.02& 140.40& 140.78& 141.16& 141.54& 141.91&  100.00 \\\hline
	 110.00& 142.29& 142.67& 143.05& 143.43& 143.80& 144.18& 144.56& 144.94& 145.31& 145.69&  110.00 \\\hline
	 120.00& 146.07& 146.44& 146.82& 147.20& 147.57& 147.95& 148.33& 148.70& 149.08& 149.46&  120.00 \\\hline
	 130.00& 149.83& 150.21& 150.58& 150.96& 151.33& 151.71& 152.08& 152.46& 152.83& 153.21&  130.00 \\\hline
	 140.00& 153.58& 153.96& 154.33& 154.71& 155.08& 155.46& 155.83& 156.20& 156.58& 156.95&  140.00 \\\hline
	 150.00& 157.33& 157.70& 158.07& 158.45& 158.82& 159.19& 159.56& 159.94& 160.31& 160.68&  150.00 \\\hline
	 160.00& 161.05& 161.43& 161.80& 162.17& 162.54& 162.91& 163.29& 163.66& 164.03& 164.40&  160.00 \\\hline
	 170.00& 164.77& 165.14& 165.51& 165.89& 166.26& 166.63& 167.00& 167.37& 167.74& 168.11&  170.00 \\\hline 
	 180.00& 168.48& 168.85& 169.22& 169.59& 169.96& 170.33& 170.70& 171.07& 171.43& 171.80&  180.00 \\\hline
	 190.00& 172.17& 172.54& 172.91& 173.28& 173.65& 174.02& 174.38& 174.75& 175.12& 175.49&  190.00 \\\hline
\hline
	\textbf{\degres C} & \textbf{0} & \textbf{1} & \textbf{2} & \textbf{3} & \textbf{4} & \textbf{5} & \textbf{6} & \textbf{7} & \textbf{8} & \textbf{9} & \textbf{\degres C}\\\hline\hline
	 200.00& 175.86& 176.22& 176.59& 176.96& 177.33& 177.69& 178.06& 178.43& 178.79& 179.16&  200.00 \\\hline
	 210.00& 179.53& 179.89& 180.26& 180.63& 180.99& 181.36& 181.72& 182.09& 182.46& 182.82&  210.00 \\\hline
	 220.00& 183.19& 183.55& 183.92& 184.28& 184.65& 185.01& 185.38& 185.74& 186.11& 186.47&  220.00 \\\hline
	 230.00& 186.84& 187.20& 187.56& 187.93& 188.29& 188.66& 189.02& 189.38& 189.75& 190.11&  230.00 \\\hline
	 240.00& 190.47& 190.84& 191.20& 191.56& 191.92& 192.29& 192.65& 193.01& 193.37& 193.74&  240.00 \\\hline
	 250.00& 194.10& 194.46& 194.82& 195.18& 195.55& 195.91& 196.27& 196.63& 196.99& 197.35&  250.00 \\\hline
	 260.00& 197.71& 198.07& 198.43& 198.79& 199.15& 199.51& 199.87& 200.23& 200.59& 200.95&  260.00 \\\hline
	 270.00& 201.31& 201.67& 202.03& 202.39& 202.75& 203.11& 203.47& 203.83& 204.19& 204.55&  270.00 \\\hline
	 280.00& 204.90& 205.26& 205.62& 205.98& 206.34& 206.70& 207.05& 207.41& 207.77& 208.13&  280.00 \\\hline
	 290.00& 208.48& 208.84& 209.20& 209.56& 209.91& 210.27& 210.63& 210.98& 211.34& 211.70&  290.00 \\\hline
	 300.00& 212.05& 212.41& 212.76& 213.12& 213.48& 213.83& 214.19& 214.54& 214.90& 215.25&  300.00 \\\hline
	 310.00& 215.61& 215.96& 216.32& 216.67& 217.03& 217.38& 217.74& 218.09& 218.44& 218.80&  310.00 \\\hline
	 320.00& 219.15& 219.51& 219.86& 220.21& 220.57& 220.92& 221.27& 221.63& 221.98& 222.33&  320.00 \\\hline 
  	 330.00& 222.68& 223.04& 223.39& 223.74& 224.09& 224.45& 224.80& 225.15& 225.50& 225.85& 330.00 \\\hline
	 340.00& 226.21& 226.56& 226.91& 227.26& 227.61& 227.96& 228.31& 228.66& 229.02& 229.37& 340.00 \\\hline
	 350.00& 229.72& 230.07& 230.42& 230.77& 231.12& 231.47& 231.82& 232.17& 232.52& 232.87& 350.00 \\\hline
	 360.00& 233.21& 233.56& 233.91& 234.26& 234.61& 234.96& 235.31& 235.66& 236.00& 236.35& 360.00 \\\hline
	 370.00& 236.70& 237.05& 237.40& 237.74& 238.09& 238.44& 238.79& 239.13& 239.48& 239.83& 370.00 \\\hline
	 380.00& 240.18& 240.52& 240.87& 241.22& 241.56& 241.91& 242.26& 242.60& 242.95& 243.29& 380.00 \\\hline
	 390.00& 243.64& 243.99& 244.33& 244.68& 245.02& 245.37& 245.71& 246.06& 246.40& 246.75& 390.00 \\\hline
\hline
	\textbf{\degres C} & \textbf{0} & \textbf{1} & \textbf{2} & \textbf{3} & \textbf{4} & \textbf{5} & \textbf{6} & \textbf{7} & \textbf{8} & \textbf{9} & \textbf{\degres C}\\\hline\hline
	 400.00& 247.09& 247.44& 247.78& 248.13& 248.47& 248.81& 249.16& 249.50& 249.85& 250.19& 400.00 \\\hline
	 410.00& 250.53& 250.88& 251.22& 251.56& 251.91& 252.25& 252.59& 252.93& 253.28& 253.62& 410.00 \\\hline
	 420.00& 253.96& 254.30& 254.65& 254.99& 255.33& 255.67& 256.01& 256.35& 256.70& 257.04& 420.00 \\\hline
	 430.00& 257.38& 257.72& 258.06& 258.40& 258.74& 259.08& 259.42& 259.76& 260.10& 260.44& 430.00 \\\hline
	 440.00& 260.78& 261.12& 261.46& 261.80& 262.14& 262.48& 262.82& 263.16& 263.50& 263.84& 440.00 \\\hline
	 450.00& 264.18& 264.52& 264.86& 265.20& 265.53& 265.87& 266.21& 266.55& 266.89& 267.22& 450.00 \\\hline
	 460.00& 267.56& 267.90& 268.24& 268.57& 268.91& 269.25& 269.59& 269.92& 270.26& 270.60& 460.00 \\\hline
	 470.00& 270.93& 271.27& 271.61& 271.94& 272.28& 272.61& 272.95& 273.29& 273.62& 273.96& 470.00 \\\hline
	 480.00& 274.29& 274.63& 274.96& 275.30& 275.63& 275.97& 276.30& 276.64& 276.97& 277.31& 480.00 \\\hline
	 490.00& 277.64& 277.98& 278.31& 278.64& 278.98& 279.31& 279.64& 279.98& 280.31& 280.64& 490.00 \\\hline
	 500.00& 280.98& 281.31& 281.64& 281.98& 282.31& 282.64& 282.97& 283.31& 283.64& 283.97& 500.00 \\\hline
	 510.00& 284.30& 284.63& 284.97& 285.30& 285.63& 285.96& 286.29& 286.62& 286.95& 287.29& 510.00 \\\hline
	 520.00& 287.62& 287.95& 288.28& 288.61& 288.94& 289.27& 289.60& 289.93& 290.26& 290.59& 520.00 \\\hline
	 530.00& 290.92& 291.25& 291.58& 291.91& 292.24& 292.56& 292.89& 293.22& 293.55& 293.88& 530.00 \\\hline
	 540.00& 294.21& 294.54& 294.86& 295.19& 295.52& 295.85& 296.18& 296.50& 296.83& 297.16& 540.00 \\\hline
	 550.00& 297.49& 297.81& 298.14& 298.47& 298.80& 299.12& 299.45& 299.78& 300.10& 300.43& 550.00 \\\hline
	 560.00& 300.75& 301.08& 301.41& 301.73& 302.06& 302.38& 302.71& 303.03& 303.36& 303.69& 560.00 \\\hline
	 570.00& 304.01& 304.34& 304.66& 304.98& 305.31& 305.63& 305.96& 306.28& 306.61& 306.93& 570.00 \\\hline
	 580.00& 307.25& 307.58& 307.90& 308.23& 308.55& 308.87& 309.20& 309.52& 309.84& 310.16& 580.00 \\\hline
	 590.00& 310.49& 310.81& 311.13& 311.45& 311.78& 312.10& 312.42& 312.74& 313.06& 313.39& 590.00 \\\hline
\hline
	\textbf{\degres C} & \textbf{0} & \textbf{1} & \textbf{2} & \textbf{3} & \textbf{4} & \textbf{5} & \textbf{6} & \textbf{7} & \textbf{8} & \textbf{9} & \textbf{\degres C}\\\hline\hline
	 600.00& 313.71& 314.03& 314.35& 314.67& 314.99& 315.31& 315.64& 315.96& 316.28& 316.60& 600.00 \\\hline
	 610.00& 316.92& 317.24& 317.56& 317.88& 318.20& 318.52& 318.84& 319.16& 319.48& 319.80& 610.00 \\\hline
	 620.00& 320.12& 320.43& 320.75& 321.07& 321.39& 321.71& 322.03& 322.35& 322.67& 322.98& 620.00 \\\hline
	 630.00& 323.30& 323.62& 323.94& 324.26& 324.57& 324.89& 325.21& 325.53& 325.84& 326.16& 630.00 \\\hline
	 640.00& 326.48& 326.79& 327.11& 327.43& 327.74& 328.06& 328.38& 328.69& 329.01& 329.32& 640.00 \\\hline
	 650.00& 329.64& 329.96& 330.27& 330.59& 330.90& 331.22& 331.53& 331.85& 332.16& 332.48& 650.00 \\\hline
	 660.00& 332.79& 333.11& 333.42& 333.74& 334.05& 334.36& 334.68& 334.99& 335.31& 335.62& 660.00 \\\hline
	 670.00& 335.93& 336.25& 336.56& 336.87& 337.18& 337.50& 337.81& 338.12& 338.44& 338.75& 670.00 \\\hline
	 680.00& 339.06& 339.37& 339.69& 340.00& 340.31& 340.62& 340.93& 341.24& 341.56& 341.87& 680.00 \\\hline
	 690.00& 342.18& 342.49& 342.80& 343.11& 343.42& 343.73& 344.04& 344.35& 344.66& 344.97& 690.00 \\\hline
	 700.00& 345.28& 345.59& 345.90& 346.21& 346.52& 346.83& 347.14& 347.45& 347.76& 348.07& 700.00 \\\hline
	 710.00& 348.38& 348.69& 348.99& 349.30& 349.61& 349.92& 350.23& 350.54& 350.84& 351.15& 710.00 \\\hline
	 720.00& 351.46& 351.77& 352.08& 352.38& 352.69& 353.00& 353.30& 353.61& 353.92& 354.22& 720.00 \\\hline
	 730.00& 354.53& 354.84& 355.14& 355.45& 355.76& 356.06& 356.37& 356.67& 356.98& 357.28& 730.00 \\\hline
	 740.00& 357.59& 357.90& 358.20& 358.51& 358.81& 359.12& 359.42& 359.72& 360.03& 360.33& 740.00 \\\hline
	 750.00& 360.64& 360.94& 361.25& 361.55& 361.85& 362.16& 362.46& 362.76& 363.07& 363.37& 750.00 \\\hline
	 760.00& 363.67& 363.98& 364.28& 364.58& 364.89& 365.19& 365.49& 365.79& 366.10& 366.40& 760.00 \\\hline
	 770.00& 366.70& 367.00& 367.30& 367.60& 367.91& 368.21& 368.51& 368.81& 369.11& 369.41& 770.00 \\\hline
	 780.00& 369.71& 370.01& 370.31& 370.61& 370.91& 371.21& 371.51& 371.81& 372.11& 372.41& 780.00 \\\hline
	 790.00& 372.71& 373.01& 373.31& 373.61& 373.91& 374.21& 374.51& 374.81& 375.11& 375.41& 790.00 \\\hline
	 800.00& 375.70& 376.00& 376.30& 376.60& 376.90& 377.19& 377.49& 377.79& 378.09& 378.39& 800.00 \\\hline
	 810.00& 378.68& 378.98& 379.28& 379.57& 379.87& 380.17& 380.46& 380.76& 381.06& 381.35& 810.00 \\\hline
	 820.00& 381.65& 381.95& 382.24& 382.54& 382.83& 383.13& 383.42& 383.72& 384.01& 384.31& 820.00 \\\hline
	 830.00& 384.60& 384.90& 385.19& 385.49& 385.78& 386.08& 386.37& 386.67& 386.96& 387.25& 830.00 \\\hline
	 840.00& 387.55& 387.84& 388.14& 388.43& 388.72& 389.02& 389.31& 389.60& 389.90& 390.19& 840.00 \\\hline
	 850.00& 390.48&       &       &       &       &       &       &       &       &       & 850.00 \\\hline
\hline
	\textbf{\degres C} & \textbf{0} & \textbf{1} & \textbf{2} & \textbf{3} & \textbf{4} & \textbf{5} & \textbf{6} & \textbf{7} & \textbf{8} & \textbf{9} & \textbf{\degres C}\\\hline\hline

\end{supertabular}
\end{center}
\caption{Table Pt100}
}
\end{changemargin}
%\end{table}


%\thispagestyle{empty}

\chapter{Calcul d'incertitude} \label{Incertitude}
%%%
\normalsize
Ici nous allons effectuer en d�tail les calculs d'incertitude pour estimer la pr�cision de la cha�ne de mesure.
Ceci dans le but de conna�tre l'incertitudes des mesures de temp�ratures envoy�es par le PIC. Les calculs seraient all�g�s en utilisant les d�riv�es logarithmiques, mais j'ai pr�f�r� n'utiliser qu'une seule m�thode.


\section{Donn�es}
Les documentations constructeurs nous donnent~:\newline
$R_s=1~k\Omega\ � 0,1\%$;\newline
$R=100~\Omega\ � 0,1\%$;\newline
$R_{Pt100}=(R_0+\delta R)~\Omega\ $ � $0,01\%$. Pour la suite on choisira $R_0=R$ pour avoir $R_0$ � $0,1\%$ et $\delta R $ � $0,01\%$;\newline
Pour une temp�rature sup�rieure � $\theta =0~\tccentigrade$, on � $2,498~V\leq V_{ref}\leq 2,500~V$. \newline
L'OPA335 a++++ un offset d'entr�e de $5~\mu V$ et l'amplificateur d'instrumentation (INA118) un offset de $50~\mu V$;\newline
$R_g$ la r�sistance de gain, est � 0,1\%. \newline

Soit $T$ la temp�rature et $T_0$ la temp�rature de r�f�rence (0,00~\degres C dans notre cas). Pour les calculs d'incertitudes pour le cas le plus d�favorable, on choisit la temp�rature la plus �loign�e, $\theta = 30~\tccentigrade$.

\section{Les �tapes}
\subsection*{Incertitude absolue de $V_{ref}$}
$V_{ref}=\overline{V_{ref}}\pm \Delta V_{ref}$\newline
\begin{equation}
\fbox{$V_{ref}=(2,499\pm 0,001)~V$}
\end{equation}
\subsection*{Incertitude absolue sur $R_s$, $R$ et $\delta R$}
$\Delta R_s=(R_s\times 0,1\%)~\Omega$. Soit~:
\begin{equation}
\fbox{$R_s=(1000\pm 1)~\Omega$}\newline
\end{equation}
De m�me on trouve 
\begin{equation}
\fbox{$R=(100\pm 0,1)~\Omega$}\newline
\end{equation}
Pour calculer l'incertitude sur $\delta R$ on choisira le cas le plus d�favorable, c'est-�-dire la plus grande variation � mesurer. Pour $\delta R(\theta)=11,67~\Omega$. 

D'o� $\Delta \delta R=(\delta R(\theta)\times 0,01\%)=0,0012~\Omega$. Finalement, 
\begin{equation}
\fbox{$\delta R=(\delta R\pm 0,001)~\Omega$}
\end{equation}

\subsection*{Incertitude absolue sur $I$}
On a $I=\frac{V_{ref}}{R_s}$\newline
$$dI=\frac{\partial I}{\partial V_{ref}}dV_{ref} + \frac{\partial I}{\partial R_s}dR_s=\frac{1}{R_s}dV_{ref}-\frac{V_{ref}}{R_s^2}dR_s$$ \newline
On en d�duit l'incertitude relative $\Delta I$~:

$$\Delta I=\left|\frac{1}{R_s}\right|\Delta V_{ref}+\left|\frac{-V_{ref}}{R_s^2}\right|\Delta R_s$$ \newline
Application num�rique~:$\Delta I=\left|\frac{1}{100}\right|\times 0,001+\left|\frac{-2,499}{100^2}\right|\times 1$. D'o�~:
\begin{equation}
\fbox{$I=(2,499\pm 0,004)~mA$}
\end{equation}
\subsection*{Incertitude absolue sur $V_0$}
Soient $V_A$ et $V_B$ tels que $V_0=V_A-V_B$ et $I'$ le courant traversant une sonde. Il faut avant tout d�terminer $\Delta I'$.

\subsubsection{Incertitude sur $I'$}
$I'=I\frac{2R+\delta R}{4R+2\delta R}=\frac{I}{2}=1,2495$. D'o� la diff�rentielle suivante~: 
\begin{eqnarray*}
dI'&=&\frac{\partial I'}{\partial I}dI \\
\Delta I'& = & \left|\frac{1}{2}\right|\Delta I=1,7\mu A
\end{eqnarray*}

\begin{equation}
\fbox{$I'=(1,250\pm 0,002)~mA$}
\end{equation}

\subsubsection{Incertitude sur $V_A$}
On a $V_A=I'\times R$

\begin{eqnarray*}
dV_A & = &\frac{\partial V_A}{\partial I'}dI' + \frac{\partial V_A}{\partial R}dR \\
& = & RdI'+I'dR\\
\Delta V_A & = & |R|\Delta I'+|I'|\Delta R
\end{eqnarray*}

\begin{equation}
\fbox{$\Delta V_A=0,325~mV$}
\end{equation}

\subsubsection{Incertitude sur $V_B$}
On a
$$V_B = I\frac{2R+\delta R}{4R+2\delta R}(2R+\delta R)=I\frac{(2R+\delta R)^2}{4R+2\delta R}$$

\begin{eqnarray*}
dV_B & = &\frac{\partial V_B}{\partial I}dI + \frac{\partial V_B}{\partial R}dR + \frac{\partial V_B}{\partial \delta R}d\delta R\\
& = &\frac{(2R+\delta R)^2}{4R+2\delta R}dI +I\frac{(8R+4\delta R)(4R+2\delta R)-4(2R+\delta R)^2}{(4R+2\delta R)^2}dR +\\
& &  I\frac{(4R+2\delta R)(4R+2\delta R)-2(2R+\delta R)^2}{(4R+2\delta R)^2}d\delta R\\
& = & \frac{(2R+\delta R)^2}{4R+2\delta R}dI +I\frac{16R^2-8\delta R^2}{(4R+2\delta R)^2}dR + I\frac{1}{2}d\delta R
\end{eqnarray*}
Pour un cas g�n�ral on obtient donc~:
\begin{equation}
\fbox{$\Delta V_B=\left|\frac{(2R+\delta R)^2}{4R +\delta R}\right|\Delta I+\left|I\frac{16R^2-8\delta R^2}{(4R+2\delta R)^2}\right|\Delta R+\left|\frac{1}{2}\right|\Delta \delta R$}
\end{equation}

Dans notre �tude, nous allons prendre le cas le plus d�favorable, c'est-�-dire $T=\theta$. Soit $\delta R=11,67~\Omega$.\newline
Application num�rique~:
\begin{equation}
\fbox{$\Delta V_B(\delta R(\theta))=1,157~mV$}
\end{equation}

\subsubsection{Incertitude sur $V_0$}
En reprenant le calcul r�alis� dans le rapport (page \pageref{V0}) on a~:$$V_0=\frac{I}{2}\times \delta R$$
Et pour l'incertitude, 
$$\Delta V_0=\Delta V_A+\Delta V_B=1,482~mV$$
D'o�
\begin{equation}
\fbox{$ V_0=(\frac{I}{2}\times \delta R\pm 1,482\times 10^{-3})~V$}
\end{equation}


\subsection*{Incertitude absolue sur $V_s$}
\subsubsection{Gain}
Soit $V_s$ le signal disponible en sortie de l'amplificateur d'instrumentation, et $R_g$ sa r�sistance de gain � 0,1\%. La d�termination de $R_g$ est issue de la formule disponible dans la documentation constructeur (voir ci-dessous) pour une pleine �chelle (0~V--5~V). 
\begin{equation}
\fbox{$R_g=(150\pm 0,15)~\Omega$}\newline
\end{equation}

On a 
\begin{eqnarray*}
G & = & 1+\frac{50\times 10^3~[\Omega]}{R_g}=334,33\\
dG & = & \frac{\partial G}{\partial R_g}dR_g = \frac{-200\times 10^3}{R_g^2}dR_g\\
\Delta G & = &  \left|\frac{50\times 10^3}{R_g^2}\right|\Delta R_g= 0,3333~\Omega 
\end{eqnarray*}

Soit
\begin{equation}
\fbox{$G=334,33\pm 0,3333$}\newline
\end{equation}

\subsubsection{Calcul de l'incertitude sur $V_s$}
\begin{eqnarray*}
V_s & = & V_0\times G\\
dV_s & = & \frac{\partial V_s}{\partial V_0}dV_0 + \frac{\partial V_s}{\partial G}dG=GdV_0+V_0dG\\
\Delta V_s(\theta) & = & G\cdot\Delta V_0 + V_0\cdot\Delta G
\end{eqnarray*}

Application num�rique pour $T=\theta$~:$\Delta V_s=0,5003~V$. On en d�duit donc que les $5~\mu$V d'offset de l'INA118 sont n�gligeables.

\begin{equation}
\fbox{$ V_s=(G\cdot\frac{I}{2}\cdot\delta R\pm 0,5003)~V$}
\end{equation}

\subsection*{Calcul final}
Des calculs pr�c�dents, on trouve (avec l'incertitude valable pour une plage 0--30~\degres C)~:
\begin{equation}
V_s=\left (G\cdot\frac{I}{2}\cdot\delta R \pm 0,5003\right )~V
\label{Vs}
\end{equation}

Les tables (annexe~\ref{tablePt100}) donnent la r�sistance standard $R_{Pt100}$ de la Pt100 � la temp�rature $T$, pour une r�sistance nominale $R_0=100,00$ Ohm � la temp�rature de r�f�rence $T_0=0~\tccentigrade $. D�pendance pouvant �tre approxim�e par la relation~:

$R_{Pt100}=R_0\left (1+A\cdot\Delta T +B\cdot\Delta T ^2\right )$, o�~: $\Delta T =T -T _0$. Avec $A=3,9083\times 10^{-3}~K^{-1}$ et $B=-5,775\times 10^{-7}~K^{-2}$.
Ce qui donne pour d�terminer la temp�rature~:
$T=T_0+\frac{\sqrt{A^2+4B\cdot r}-A}{2B}$, o� 
\begin{equation}
r=\frac{R_{Pt100}}{R_0}-1
\label{r}
\end{equation}


Or on a $R_0=R$ et $R_{Pt100}=R+\delta R$. Notre variable �tant $V_s$, la tension disponible en entr�e du PIC, il ne faut obtenir que $V_s$ pour calculer la temp�rature. On l'obtient par la relation \eqref{Vs}. Soit,
$$ \delta R=\frac{2V_s}{G\cdot I}$$

d'o�~:
	
 

\begin{eqnarray*}
r & = & \frac{R+ \delta R}{R}-1 = \frac{\delta R}{R} = \frac{2V_s}{G\cdot I\cdot R}\\
dr & = & \frac{\partial r}{\partial V_s}dV_s + \frac{\partial r}{\partial G}dG+ \frac{\partial r}{\partial I}dI+ \frac{\partial r}{\partial R}dR\\
& = & \frac{2}{G\cdot I\cdot R}dV_s-2V_s\left ( \frac{dG}{G^2\cdot I\cdot R} + \frac{dI}{G\cdot I^2\cdot R} + \frac{dR}{G\cdot I\cdot R^2} \right )\\
& = & \frac{2}{G\cdot I\cdot R}dV_s-2\frac{G\cdot I\cdot\delta R}{2}\left ( \frac{dG}{G^2\cdot I\cdot R} + \frac{dI}{G\cdot I^2\cdot R} + \frac{dR}{G\cdot I\cdot R^2} \right )\\
\Delta r & = & \frac{2}{G\cdot I\cdot R}\Delta V_s+\delta R\left ( \frac{\Delta G}{G\cdot R} + \frac{\Delta I}{I\cdot R} + \frac{\Delta R}{R^2} \right )
\end{eqnarray*}

Soit l'application num�rique pour le cas le plus d�favorable~: $\Delta r = 0,0124$


Finalement on obtient~:
\begin{equation}
\fbox{$T=T_0+\frac{\sqrt{A^2+4B\cdot r}-A}{2B}$}
\label{Tfin}
\end{equation}
Avec
\begin{equation}
\fbox{$r=\frac{2V_s}{G\cdot I\cdot R}\pm 0,0124 $}
\label{Rfin}
\end{equation}


\subsection*{Incertitude finale}
On reprend $r  =  \frac{\delta R(\theta)}{R}$ pour calculer l'incertitude sur $T$~:
\begin{eqnarray*}
T & = & T_0+\frac{\sqrt{A^2+4B\cdot r}-A}{2B}\\
dT & = & \frac{\partial T}{\partial r}dr = \frac{1}{2B}\cdot\frac{4B}{2\sqrt{A^2+4B\cdot r}}dr = \frac{1}{\sqrt{A^2+4B\cdot r}}dr\\
\Delta T & = & \left |\frac{1}{\sqrt{A^2+4B\cdot r}}\right | \Delta r
\end{eqnarray*}


Application num�rique~:$\Delta T = 0,066~\tccentigrade$


\subsection*{R�sultat}

Th�oriquement, on obtient une mesure fid�le � 0,1~\degres C (avec une marge de 0,03~\degres C). Avec pour relations~:
\begin{center}
	\shadowbox{
\vspace{4mm}
		\begin{Beqnarray} 
		T & = & \left (T_0+\frac{\sqrt{A^2+4B\cdot r}-A}{2B}\pm 0,07\right ) \tccentigrade \label{rslt}\\
		r & = & \frac{2V_s}{G\cdot I\cdot R}\pm 0,0124 \label{rslt2}\\
		A & = & 3,9083\cdot 10^{-3}~K^{-1} \\
		B & = & -5,775\cdot 10^{-7}~K^{-2}\\
		T_0 & = & 0,00~\tccentigrade
		\end{Beqnarray}
\vspace{4mm}
	}
\end{center}

Avec pour le circuit pr�sent�~:
\begin{eqnarray}
	G & = & 425,53\pm 0,4255 \\
	I & = & (2,499\pm 0,004)~mA \\
	R & = & (100\pm 0,1)~\Omega 
\end{eqnarray}


%\thispagestyle{empty}

\chapter{Cartes}

\section{Emetteur}
Vous trouverez page suivante le sch�ma de la carte �mettrice. Celle qui mesure la temp�rature. Pour un soucis de consommation, plusieurs r�gulateurs avec "Enable" ont �t� utilis�s pour pouvoir couper l'alimentation de tout ce qui est inutile � certain moment.

Puis les pages suivantes sont directement tir�e du fichier PostScript export� par Pcb. Il s'agit des typons de la carte �mettrice. 

\section{R�cepteur}
Ensuite vous trouverez le sch�ma de la carte r�cepteur. Bien remarquer le condensateur de 470~nF. C'est par lui qu'est fix� le potentiel 3,3~V pour fixer la vitesse de transmission sur la liaison USB. 

Puis les typons de la carte r�cepteur export�s par Pcb. 3 feuilles parmis 12 (masque, trou de per�age, \dots)

\begin{landscape}
%%\addcontentsline{toc}{chapter}{Sch�ma}
\begin{figure}[h]
    \includegraphics[]{./img/Emet}
    \caption{Schema \' Emetteur}
\end{figure}
\end{landscape}

\includepdf[pages={2,3,11},pagecommand={\thispagestyle{plain}}]{Emetteur} 



\begin{landscape}
\begin{figure}[]
    \includegraphics[width=25cm]{./img/Recep}
    \caption{Schema \' Emetteur}
\end{figure}
\end{landscape}
\includepdf[pages={2,3,11},pagecommand={\thispagestyle{plain}}]{Recepteur} 


%\includepdf[pages={{},1},nup=1x2,landscape=true]{output}
%\includepdf[pages={{},2},nup=1x2,landscape=true,angle=180]{output}
%\includepdf[pages={6,3},nup=1x2,landscape=true]{output}
%\includepdf[pages={5,4},nup=1x2,landscape=true,angle=180]{output} 


%\chapter{Organigramme}
%\addcontentsline{toc}{chapter}{D~~Organigramme}
%\includepdf[pages={1},nup=1x1,landscape=true]{AB_CO} %% nup pour l'echelle

%\includepdf[pages=-]{./source/AB_CO}

%\addtocontents{toc}{\vskip 30 mm \Large \centerline {\textsc{Article Connexe}}  \vskip 10 mm  
%\normalsize Utilisation de gEDA pour la Conception �lectronique Assist�e par Ordinateur. Utilisation de Piklab pour programmer les PIC sous Linux avec l'ICD2. Utilisation de la carte Pt100 et Xbee. Page \pageref{connexe} des annexes. \newline \vskip 5mm}

%\part*{\Huge{\textsc{Article Connexe}}}\label{connexe}
%\includepdf[pages=-]{Cern.pdf} % J'inclus un deuxieme rapport, au format pdf
%linktodoc ouvre le pdf inclu !!!! 
%% RESUME
%\setlength{\headheight}{10pt} 	% Haut de page
%\setlength{\headsep}{10pt} 	% Entre le haut de page et le texte
%\setlength{\footskip}{0pt} 	% Bas de page + s�paration
%\setlength{\textheight}{999pt} 	% Hauteur de la zone de texte 
% \begin{changemargin}{-10mm}{-5mm}

\begin{minipage}[c][150pt][c]{.40\linewidth} %[hauteur][position]{largeur}
	{\normalsize 
		\begin{flushleft}
		\vspace{3mm}
	      	CERN\\
		Organisation Europ�enne pour\\
		la Recherche Nucl�aire\\
		F-01631 CERN Cedex\\
		France\\
		\url{http://www.cern.ch/}\\
		\vspace{2mm}
		\includegraphics[width=2cm]{./img/CERN_logo}
		\vspace{3mm}
		\end{flushleft}
	}
   \end{minipage}\hfill
%}
\hspace{15mm}
%
%\Ovalbox{
   \begin{minipage}[c][150pt][c]{.40\linewidth} %[hauteur][position]{largeur}
	{\normalsize 
		\begin{flushright}
		\vspace{3mm}
		K�vin Raymond\\
		IUT Annecy\\
		D�partement GEII\\
		Ann�e 2007\\
		\vspace{2mm}
		\includegraphics[width=5cm]{./img/A}
		\vspace{3mm}
		\end{flushright}
	}
   \end{minipage}
%}
\vspace{15mm}

\begin{center}
\textbf{\large{R�sum�}}
\end{center}

\thispagestyle{empty}
\begin{center}
\textbf{Mesure d'une temp�rature avec pr�cision et transmission par liaison sans fil}
\end{center}




%\fontsize{11pt}{1}\selectfont

Dans cette �tude, nous abordons la mesure d'une temp�rature ambiante 
avec pr�cision. Apr�s avoir d�crit le fonctionnement globale d'une 
telle mesure et des diff�rentes possibilit�s de mises en \oe uvre, 
nous �tudierons le d�faut de la cha�ne de mesure et l'incertitude introduite.

Ceci passe par l'utilisation des logiciels libres de Conception �lectronique 
Assist� par Ordinateur et de programation des microcontr�leurs, ainsi que 
par l'utilisation de divers protocoles comme le ZigBee ou l'USB.

Pour une plus grande portabilit�, l'�tude est r�alis�e dans un environnement GNU/Linux.
\vspace{0.25cm}

\noindent\textbf{Mots-cl�s~:} Mesure de temp�rature, RTD, PRTD, Pt100, ZigBee, Xbee, USB, PIC, GNU, Linux. %\fontsize{12pt}{1.5}\selectfont


\vspace{10mm}
\begin{center}
	\underline{\hspace{3cm}} % Trace une fine ligne de 3 cm...
\end{center}
\vspace{10mm}
%\vskip15mm
\selectlanguage{english}
\begin{center}
\textbf{\large{Abstract}}
\end{center}
\begin{center}
\textbf{Mesure of high precision temperature and wireless equipment}
\end{center}

%\vskip5mm



%\fontsize{10pt}{1}\selectfont

In this study, we move on the precision sensor temperature. After a global description of some possibilities
of exploitation, we will study the default of the circuit conditioning.

Getting started with some open source software like PIC programmation's and Electronic Design Automation tools.
With some new protocoles like USB and Zigbee solutions.

Of course, this study is based with a GNU/Linux computeur.


\vspace{0.25cm}
\noindent\textbf{Keywords~:} Temperature sensor, RTD, PRTD, Pt100, ZigBee, Xbee, USB, PIC, GNU, Linux.
\fontsize{12pt}{1.5}\selectfont


\vspace{13mm}
\selectlanguage{french}
\begin{center}
		Rapport disponible~: \url{http://kevin.raymond.free.fr/Stage/rapport/RapportKR.pdf}

\end{center}
\end{changemargin}



\end{document}
